\documentclass[11pt]{article}
\usepackage[utf8]{inputenc}
\usepackage[T1]{fontenc}
\usepackage{booktabs}
\usepackage{amsmath}
\usepackage{geometry}
\geometry{margin=2.5cm}

\title{Causal Fairness Analysis Report}
\author{Generated by FairMind and an LLM}
\date{2026-08-22}

\begin{document}
\maketitle

\section{Analysis setup}

\begin{tabular}{ll}
\toprule
Dataset & oulad\_final\_result \\
Protected attribute ($X$) & S1\_gender \\
Comparison & F $\to$ M \\
Outcome ($Y$) & T\_final\_result (Pass) \\
Mediator ($W$) & studied\_credits \\
Confounder ($Z$) & highest\_education \\
\bottomrule
\end{tabular}

\section{Causal effects (FairMind, exact values)}

The five effects below are computed by FairMind through exact inference on the
fitted Bayesian network. These values are final and must not be recomputed or
altered.

\begin{tabular}{lr}
\toprule
Effect & Value \\
\midrule
Total Variation (TV) & -0.027800 \\
Total Effect (TE) & -0.036400 \\
Spurious Effect (SE) & 0.008600 \\
Direct Effect (DE) & -0.024400 \\
Indirect Effect (IE, decomposition form: TE = DE + IE) & -0.012000 \\
\bottomrule
\end{tabular}

\section{Qualitative interpretation}

\subsection{Total effect and spurious component (TV, TE, SE)}

The total disparity between the two groups is -0.0278, indicating a slight advantage for the protected group. After removing the confounding variable, the genuine causal effect (TE) is -0.0364, which is larger in magnitude than the total variation, suggesting that the disparity is predominantly due to genuine causal factors rather than spurious ones. The spurious effect (SE) of 0.0086 is relatively small, indicating that the confounder has a minor impact on the observed disparity.

\subsection{Direct effect (DE)}

The direct effect (DE) is -0.0244, indicating a direct disadvantage for the protected group. This suggests there is direct discrimination against the protected group, though the magnitude is not extremely large compared to the total effect.

\subsection{Indirect effect (IE)}

The indirect effect (IE) is -0.0120, which is negative and has the same sign as the direct effect (DE). This means the indirect effect through the mediator amplifies the direct effect, making the total effect more pronounced. The indirect effect is not negligible, as it contributes significantly to the overall disparity.

\section{Recap Questions}

\begin{enumerate}
\item Is there direct discrimination against the protected group? \textbf{Answer:} YES
\item Does the indirect effect through the mediator mitigate the total effect? \textbf{Answer:} NO
\item Does the observed total effect exceed the practical relevance threshold? \textbf{Answer:} NO
\item Does the spurious component (SE) contribute substantially to the observed difference? \textbf{Answer:} NO
\item Is the direct effect the dominant component compared to the indirect effect? \textbf{Answer:} YES
\end{enumerate}

\vspace{1em}
\noindent\footnotesize
The recap answers follow deterministic threshold rules applied to the values above: direct discrimination when $|\mathrm{DE}| > 0.05$; a mediator channel counted as negligible when $|\mathrm{IE}| < 0.005$; practical relevance when $|\mathrm{TV}| > 0.05$; a substantial spurious component when $|\mathrm{SE}| / |\mathrm{TV}| > 0.1$.

\end{document}